\documentclass[10pt,a4paper]{article}
\usepackage{webclases-lesson}
\lessonSetup{T2}{v0.2}{T2 - Movimiento en el plano}

\begin{document}
\shorthandoff{<>}
\color{Ink}

\coverHeader{T2 --- Movimiento en el plano}{Componentes, proyectiles y movimiento circular con un metodo comun.}

\begin{lessoncard}{Proposito}
En dos dimensiones no aparecen leyes nuevas de cinematica: aparece la necesidad de separar componentes y volver a unirlas con el mismo tiempo. La leccion agrupa los tiros para que no parezcan problemas distintos cada vez que cambia el angulo.
\begin{keybox}
\textbf{Clave de lectura.} En proyectiles ideales, $x$ e $y$ comparten el mismo reloj. En circular, rapidez constante no significa velocidad constante.
\end{keybox}
\end{lessoncard}

\begin{lessoncard}{Indice por bloques}
\footnotesize
\begin{tabularx}{\linewidth}{@{}>{\raggedright\arraybackslash}p{1.6cm}Y>{\raggedleft\arraybackslash\bfseries}p{1.0cm}@{}}
\textbf{Bloque} & \textbf{Alcance} & \textbf{Pag.}\\
\midrule
\tocBlock{Bloque 1}{Movimiento en dos dimensiones}
\tocLine{sec:b1}{1}{Vectores, componentes, velocidad tangente y velocidad relativa}
\tocBlock{Bloque 2}{Movimiento de proyectiles}
\tocLine{sec:b2}{2}{Modelo, tiros horizontal y oblicuo, alturas distintas y graficas}
\tocBlock{Bloque 3}{Movimiento circular}
\tocLine{sec:b3}{3}{Magnitudes angulares, aceleraciones, MCU y MCUA}
\tocBlock{Bloque 4}{Elegir modelo y comprobar}
\tocLine{sec:b4}{4}{Checklist de decision y errores frecuentes}
\end{tabularx}
\end{lessoncard}

{\scriptsize\textcolor{Muted}{Fuentes: PDF original T2, background editorial T2 y contraste conceptual con OpenStax University Physics 4. Diagramas redibujados.}}

\clearpage

\lessonHeader{Bloque 1 \textperiodcentered\ Movimiento en dos dimensiones}{Apartados internos}{Componentes y trayectoria}{Objetivo: separar vectores sin perder la interpretacion geometrica.}
\begin{lessoncard}{Vectores de movimiento}
\subtopic{1}{Movimiento en dos dimensiones}{sec:b1}
\[
\vec r=x\ii+y\jj,\qquad \vec v=\frac{d\vec r}{dt}=v_x\ii+v_y\jj,\qquad \vec a=a_x\ii+a_y\jj.
\]
La velocidad instantanea es tangente a la trayectoria; la aceleracion apunta hacia donde cambia la velocidad, no necesariamente hacia donde se mueve el objeto.
\begin{center}
\begin{tikzpicture}[x=1cm,y=.8cm]
  \draw[axis] (0,0) -- (6,0) node[right] {\scriptsize $x$};
  \draw[axis] (0,0) -- (0,4) node[above] {\scriptsize $y$};
  \draw[BrandBlue,line width=1pt,domain=.4:5.2,samples=80] plot(\x,{.18*(\x-2)^2+1});
  \coordinate (P) at (3.4,1.35);
  \fill[BrandBlue] (P) circle (.08);
  \draw[vector] (P) -- ++(1.1,.28) node[right] {\scriptsize $\vec v$ tangente};
  \draw[vectoramber] (P) -- ++(.45,-.85) node[right] {\scriptsize $\vec a$ posible};
\end{tikzpicture}
\end{center}
\end{lessoncard}

\begin{examplebox}{nadador y corriente}
Un nadador avanza a $1,40\mps$ respecto al agua perpendicular a la orilla. El rio baja a $0,60\mps$ y tiene $42\units{m}$ de ancho.
\solutionblock
\[
\vec v_{\mathrm{suelo}}=0,60\ii+1,40\jj,\qquad t=\frac{42}{1,40}=30\units{s}.
\]
Deriva horizontal: $x=0,60(30)=18\units{m}$. La rapidez respecto al suelo es $\sqrt{0,60^2+1,40^2}=1,52\mps$.
\end{examplebox}

\clearpage

\lessonHeader{Bloque 2 \textperiodcentered\ Proyectiles}{Apartados internos}{Dos componentes, un tiempo}{Objetivo: resolver variantes sin duplicar el metodo.}
\begin{lessoncard}{Modelo comun}
\subtopic{2}{Movimiento de proyectiles}{sec:b2}
Hipotesis: se desprecia el rozamiento del aire y la aceleracion es $\vec a=(0,-g)$. Entonces:
\[
x=x_0+v_{0x}t,\qquad y=y_0+v_{0y}t-\frac12gt^2,\qquad v_y=v_{0y}-gt.
\]
El tiempo se obtiene normalmente desde la ecuacion vertical; luego se usa en la horizontal.
\begin{center}
\begin{tikzpicture}[x=1cm,y=.65cm]
  \draw[axis] (0,0) -- (6.1,0) node[right] {\scriptsize $x$};
  \draw[axis] (0,0) -- (0,4.2) node[above] {\scriptsize $y$};
  \draw[BrandBlue,line width=1pt,domain=0:5.6,samples=80] plot(\x,{2.9+.25*\x-.11*\x*\x});
  \draw[guide] (1,0)--(1,3.04) (3,0)--(3,2.66) (5,0)--(5,1.4);
  \foreach \x/\y in {1/3.04,3/2.66,5/1.4} \fill[BrandBlue] (\x,\y) circle (.06);
  \node[text=Muted] at (3,-.55) {\scriptsize mismo $t$ para horizontal y vertical};
\end{tikzpicture}
\end{center}
\end{lessoncard}

\begin{examplebox}{lanzamiento horizontal}
Desde una mesa de $1,25\units{m}$ se lanza una bola horizontalmente a $3,0\mps$. Usa $g=10\mpss$.
\solutionblock Vertical: $0=1,25-5t^2$, luego $t=0,50\units{s}$. Horizontal: $x=3,0(0,50)=1,5\units{m}$. Al impactar, $v_x=3,0\mps$ y $v_y=-5,0\mps$.
\variation{Otra bola soltada desde la misma altura cae en el mismo tiempo. La velocidad horizontal cambia el alcance, no el tiempo de caida ideal.}
\end{examplebox}

\clearpage

\lessonHeader{Bloque 2 \textperiodcentered\ Proyectiles}{Aplicaciones}{Oblicuo, angulos y alturas distintas}{Objetivo: reconocer que las variantes comparten ecuaciones.}
\begin{examplebox}{salto oblicuo}
Un balon sale con $v_0=12,0\mps$ y $\alpha=35^\circ$ desde el suelo. Usa $g=10\mpss$.
\solutionblock
\[
v_{0x}=12\cos35^\circ=9,83\mps,\qquad v_{0y}=12\sin35^\circ=6,88\mps.
\]
Tiempo de vuelo hasta $y=0$: $t=2v_{0y}/g=1,38\units{s}$. Alcance: $R=v_{0x}t=13,6\units{m}$. Altura maxima: $h=v_{0y}^2/(2g)=2,37\units{m}$.
\end{examplebox}

\begin{lessoncard}{Comparacion $30^\circ$ y $60^\circ$}
Con la misma rapidez inicial y misma altura de salida/llegada:
\[
R=\frac{v_0^2\sin 2\alpha}{g}.
\]
Los angulos $30^\circ$ y $60^\circ$ dan el mismo alcance porque $\sin60^\circ=\sin120^\circ$, pero no la misma altura ni el mismo tiempo.
\begin{center}
\begin{tabular}{@{}ccc@{}}
\toprule
Angulo & tiempo de vuelo & altura maxima\\
\midrule
$30^\circ$ & menor & menor\\
$60^\circ$ & mayor & mayor\\
\bottomrule
\end{tabular}
\end{center}
\end{lessoncard}

\begin{examplebox}{llegada a otra altura}
Pelota desde una ventana a $y_0=8,0\units{m}$ con $v_0=10,0\mps$ y $\alpha=20^\circ$.
\solutionblock
\[
x=9,40t,\qquad y=8,0+3,42t-5t^2.
\]
Suelo: $0=8,0+3,42t-5t^2$, raiz positiva $t=1,69\units{s}$. Alcance: $x=15,9\units{m}$.
\end{examplebox}

\clearpage

\lessonHeader{Bloque 2 \textperiodcentered\ Proyectiles}{Graficas temporales}{Que cambia y que no}{Objetivo: conectar trayectoria, $x(t)$, $y(t)$, $v_x(t)$ y $v_y(t)$.}
\begin{lessoncard}{Cuatro lecturas de un mismo tiro}
En el modelo ideal, $v_x$ es constante y $v_y$ cambia linealmente. Por eso $x(t)$ es lineal, $y(t)$ es parabolica y la trayectoria $y(x)$ tambien es parabolica.
\begin{center}
\begin{tikzpicture}[x=.75cm,y=.55cm]
  \draw[axis] (0,0)--(4.5,0) node[right] {\scriptsize $t$};
  \draw[axis] (0,0)--(0,3.2) node[above] {\scriptsize $x$};
  \draw[BrandBlue,line width=1pt] (.3,.4)--(4,2.7);
  \node at (2.4,3.0) {\scriptsize $x(t)$};
  \begin{scope}[xshift=5.3cm]
    \draw[axis] (0,0)--(4.5,0) node[right] {\scriptsize $t$};
    \draw[axis] (0,0)--(0,3.2) node[above] {\scriptsize $y$};
    \draw[BrandBlue,line width=1pt,domain=.2:4,samples=60] plot(\x,{.55+1.55*\x-.33*\x*\x});
    \node at (2.4,3.0) {\scriptsize $y(t)$};
  \end{scope}
  \begin{scope}[yshift=-4.1cm]
    \draw[axis] (0,0)--(4.5,0) node[right] {\scriptsize $t$};
    \draw[axis] (0,-1.4)--(0,2.2) node[above] {\scriptsize $v_y$};
    \draw[BrandBlue,line width=1pt] (.3,1.8)--(4,-1.0);
    \draw[guide] (0,0)--(4.3,0);
    \node at (2.5,2.0) {\scriptsize pendiente $-g$};
  \end{scope}
  \begin{scope}[xshift=5.3cm,yshift=-4.1cm]
    \draw[axis] (0,0)--(4.5,0) node[right] {\scriptsize $t$};
    \draw[axis] (0,0)--(0,2.5) node[above] {\scriptsize $v_x$};
    \draw[BrandBlue,line width=1pt] (.3,1.45)--(4,1.45);
    \node at (2.4,2.0) {\scriptsize constante};
  \end{scope}
\end{tikzpicture}
\end{center}
\end{lessoncard}

\begin{warnbox}
Si en un tiro ideal te sale $v_x$ cambiando sin rozamiento ni motor, has mezclado componentes. Si te sale un tiempo negativo para el impacto, elegiste la raiz que describe el pasado matematico, no el evento fisico.
\end{warnbox}

\clearpage

\lessonHeader{Bloque 3 \textperiodcentered\ Movimiento circular}{Apartados internos}{Angulos, radios y aceleraciones}{Objetivo: separar giro de desplazamiento lineal.}
\begin{lessoncard}{Magnitudes angulares}
\subtopic{3}{Movimiento circular}{sec:b3}
\[
\theta\ \text{en rad},\qquad \omega=\frac{d\theta}{dt},\qquad \alpha=\frac{d\omega}{dt},\qquad
s=R\theta,\qquad v=R\omega.
\]
Periodo y frecuencia: $T=1/f$ y $\omega=2\pi/T=2\pi f$.
\end{lessoncard}

\begin{lessoncard}{Tangencial y normal}
\[
a_t=R\alpha,\qquad a_n=\frac{v^2}{R}=R\omega^2.
\]
La aceleracion normal apunta hacia el centro; la tangencial cambia la rapidez. En MCU, $a_t=0$ pero $a_n\ne0$.
\begin{center}
\begin{tikzpicture}[x=1cm,y=1cm]
  \draw[BrandBlue,line width=1pt] (0,0) circle (1.4);
  \fill[BrandBlue] (1.0,1.0) circle (.08);
  \draw[vectorgreen] (1.0,1.0) -- (0.25,0.25) node[left] {\scriptsize $a_n$};
  \draw[vectoramber] (1.0,1.0) -- (1.75,.25) node[right] {\scriptsize $v$};
  \draw[guide] (0,0)--(1.0,1.0) node[midway,above] {\scriptsize $R$};
\end{tikzpicture}
\end{center}
\end{lessoncard}

\clearpage

\lessonHeader{Bloque 3 \textperiodcentered\ Movimiento circular}{Aplicaciones}{MCU y MCUA}{Objetivo: elegir variables angulares cuando simplifican.}
\begin{examplebox}{MCU}
Una rueda de radio $0,35\units{m}$ gira a $240\units{rpm}$.
\solutionblock
\[
f=240/60=4,0\units{Hz},\qquad \omega=2\pi f=25,1\units{rad/s}.
\]
\[
v=R\omega=8,80\mps,\qquad a_n=R\omega^2=221\mpss.
\]
Todos los puntos tienen la misma $\omega$, pero no la misma rapidez lineal si cambia $R$.
\end{examplebox}

\begin{examplebox}{frenado angular}
Un disco pasa de $\omega_0=20\units{rad/s}$ a reposo en $5,0\units{s}$ con aceleracion angular constante.
\solutionblock
\[
\alpha=\frac{0-20}{5,0}=-4,0\units{rad/s^2},\qquad
\Delta\theta=\omega_0t+\frac12\alpha t^2=50\units{rad}.
\]
Numero de vueltas: $50/(2\pi)=8,0$ vueltas. En un punto a $R=0,20\units{m}$, $a_t=R\alpha=-0,80\mpss$.
\end{examplebox}

\begin{warnbox}
No conviertas grados a radianes solo al final si vas a usar $s=R\theta$ o derivadas angulares. Esas relaciones usan radianes.
\end{warnbox}

\clearpage

\lessonHeader{Bloque 4 \textperiodcentered\ Elegir modelo y comprobar}{Sintesis}{La pregunta decide el camino}{Objetivo: evitar aplicar el ultimo ejemplo visto.}
\begin{lessoncard}{Decision rapida}
\subtopic{4}{Elegir modelo y comprobar}{sec:b4}
\begin{tabularx}{\linewidth}{@{}p{2.7cm}Y@{}}
\toprule
Si el enunciado habla de... & Piensa en...\\
\midrule
rio, viento, cinta, observadores & velocidades relativas: suma vectorial\\
tiro, proyectil, lanzamiento & componentes con mismo tiempo y $a_y=-g$\\
giro o rueda & variables angulares y aceleracion normal\\
frenado de giro & MCUA: $\omega=\omega_0+\alpha t$\\
\bottomrule
\end{tabularx}
\end{lessoncard}

\begin{exercisebox}{simular un proyectil}
Con el mismo $v_0=20\mps$, compara mentalmente $20^\circ$, $45^\circ$ y $70^\circ$ cuando salida y llegada estan a la misma altura. Ordena alcance, altura maxima y tiempo de vuelo.
\solutionblock Mayor alcance: $45^\circ$. Mayor altura y mayor tiempo: $70^\circ$. El de $20^\circ$ es mas bajo y breve. La simetria de alcances aparece en pares complementarios: $20^\circ$ y $70^\circ$ tendrian el mismo alcance ideal.
\end{exercisebox}

\begin{lessoncard}{Comprobaciones finales}
Un resultado de proyectil debe tener tiempo positivo, unidades correctas y una altura compatible con el signo de $v_y$. Un resultado circular debe distinguir radio, longitud de arco y angulo. Si un calculo da aceleracion normal negativa, revisa: el signo pertenece a la direccion elegida; el modulo $v^2/R$ es positivo.
\end{lessoncard}

\end{document}
